\section{Validation hors échantillon}
\label{sec:out_of_sample}

\lettrine[lines=2, findent=2pt, nindent=0pt]{L}{a} validation hors échantillon est le test le plus important de tout travail de recherche quantitative~: c'est la seule manière de distinguer un \emph{edge} structurel d'un artefact d'ajustement. Le portefeuille y obtient un Sharpe supérieur à celui de la période de conception --- résultat favorable, mais dont cette section explique aussi les limites.

\subsection{Comparaison directe}

\begin{center}
\renewcommand{\arraystretch}{1.15}
\begin{tabular}{lrrl}
\toprule
\textbf{Métrique} & \textbf{In-sample} & \textbf{Hors échantillon} & \textbf{Interprétation} \\
\midrule
Nombre de barres & 6995 & 449 & OOS $\approx$ 6\,\% de l'IS \\
Rendement total cumulé & $372.2\,\%$ & $193.3\,\%$ & Fenêtre courte, effet de composition \\
CAGR & $23.08\,\%$ & $\mathbf{181.48\,\%}$ & Annualisation d'une fenêtre de 449 barres \\
Volatilité & $37.00\,\%$ & $36.59\,\%$ & Comparable \\
Max Drawdown & $-55.00\,\%$ & $\mathbf{-17.84\,\%}$ & DD beaucoup plus contenu \\
Sharpe Ratio & $0.75$ & $\mathbf{3.02}$ & \textbf{OOS supérieur à l'IS} \\
\bottomrule
\end{tabular}
\end{center}

\begin{warningbox}
\textbf{Le CAGR hors échantillon de 181\,\% ne doit être lu ni comme une performance, ni comme une prévision.} C'est l'annualisation d'une fenêtre de 449 barres qui recouvre précisément la phase la plus favorable du marché des métaux. Étendre un rendement de quelques mois à une base annuelle amplifie mécaniquement tout ce que cette période contenait de circonstanciel. Le chiffre est rapporté pour la transparence du calcul~; il ne figure dans aucune conclusion de ce document.
\end{warningbox}

\begin{insightbox}
\textbf{Lecture critique~:} un véritable surajustement aurait typiquement produit une dégradation sévère hors période de conception~--- un Sharpe qui passe de $>1$ à proche de zéro ou même négatif. Ici le Sharpe hors échantillon est \emph{supérieur} à l'in-sample, et le repli maximal nettement plus faible. Cela réfute l'hypothèse d'un surajustement \emph{majeur}, sans pour autant établir que la performance se reproduira~: les tests de la section~\ref{sec:advanced_robustness} nuancent ce constat, et l'un d'eux échoue.
\end{insightbox}

\subsection{Zoom sur la fenêtre de validation}

\begin{figure}[htbp]
\centering
\includegraphics[width=\linewidth]{oos_zoom.png}
\caption{Zoom sur la période hors échantillon (courbe orange), base 100 à la date de coupure. La courbe in-sample est conservée en arrière-plan pour référence.}
\end{figure}

\subsection{Caveat méthodologique sur le caractère OOS}

\begin{warningbox}
La fenêtre hors échantillon est un découpage \emph{ex-post}~--- elle a été isolée après coup, lorsque ces données sont devenues disponibles. Il ne s'agit donc pas d'un \emph{holdout} pré-engagé avant le design, et elle recouvre en partie la tranche que la politique de \emph{holdout} du projet a gelée pour la sélection de modèle. Elle est ici utilisée pour \emph{rapporter}, jamais pour choisir un paramètre.

Trois éléments soutiennent néanmoins le résultat~:
\begin{enumerate}
\item Le Sharpe hors échantillon est \emph{supérieur} à l'in-sample, pas simplement \emph{proche}~--- c'est l'opposé d'une dégradation par surajustement.
\item Le repli maximal hors échantillon est bien plus faible qu'en in-sample, ce qui n'est pas non plus cohérent avec un surajustement.
\item Les hyperparamètres de chaque moteur sont restés inchangés lors des évolutions successives de l'allocation.
\end{enumerate}
\textbf{Action corrective retenue~:} un \emph{paper-trade} d'un minimum de trois mois avant tout déploiement de capital réel, pour consolider la validation par une fenêtre authentiquement pré-engagée.
\end{warningbox}
